\chapter{Introduzione e Descrizione del Sistema}
\label{cap:intro}

\section{Panoramica dell'Applicazione}
Il progetto oggetto di questa relazione, denominato \textbf{Guided Tours}, consiste nella progettazione e nello sviluppo di un'applicazione web orientata alla gestione di visite guidate. Il sistema è stato concepito per rispondere alle esigenze di organizzazioni locali senza scopo di lucro che si occupano della valorizzazione del patrimonio territoriale, sia esso di natura storico-artistica che paesaggistica.

L'obiettivo primario della piattaforma è la digitalizzazione del flusso organizzativo che intercorre tra la domanda, rappresentata da visitatori singoli o gruppi (scuole, turisti), e l'offerta, costituita da una rete di guide volontarie. Attualmente, tali processi sono spesso gestiti in modo frammentato e manuale; il sistema proposto mira a centralizzare le informazioni, ottimizzare l'assegnazione delle risorse umane e garantire un canale di comunicazione strutturato.

L'architettura logica dell'applicazione prevede tre attori principali, ciascuno caratterizzato da specifici privilegi e interfacce dedicate:
\begin{itemize}
    \item \textbf{Fruitore:} L'utente esterno che naviga il catalogo delle offerte culturali e inoltra richieste di prenotazione. Può agire come singolo o come capogruppo.
    \item \textbf{Volontario:} La risorsa operativa dell'organizzazione. Accede al sistema per comunicare le proprie disponibilità temporali e le proprie competenze (es. lingue parlate, abilitazioni per specifici luoghi).
    \item \textbf{Configuratore:} Il ruolo amministrativo che detiene il controllo completo del sistema. Si occupa del popolamento delle anagrafiche (Luoghi, Tipologie di Visita) e della gestione operativa delle richieste, incrociando le esigenze dei fruitori con le disponibilità dei volontari.
\end{itemize}

\section{Analisi dell'Interfaccia e Funzionalità}
L'interfaccia utente è stata progettata seguendo i principi di usabilità e chiarezza informativa, differenziando nettamente le aree pubbliche da quelle riservate. Di seguito viene fornita un'analisi dettagliata delle sezioni costitutive del sistema.

\subsection{Landing Page e Navigazione Pubblica}
La Home Page funge da ambiente di accoglienza e snodo direzionale primario per le diverse tipologie di utenza.
Graficamente, il layout è dominato da una sezione introduttiva (\textit{Hero Section}) concepita per comunicare immediatamente la missione dell'organizzazione e instaurare un rapporto di fiducia.
La barra di navigazione superiore è un elemento persistente e offre un \textit{affordance} chiara per l'accesso rapido alle funzioni di autenticazione.

Dal punto di vista dell'analisi dei task, questa vista deve assolvere simultaneamente a due obiettivi:
\begin{enumerate}
    \item \textbf{Informativo (Discovery):} Presentare la proposta di valore ai visitatori anonimi, invitandoli alla conversione (es. registrazione) al fine di prenotare una visita.
    \item \textbf{Direzionale (Routing):} Fornire un accesso immediato e privo di frizioni all'area riservata per lo staff e i volontari già registrati, agendo da \textit{gateway}.
\end{enumerate}

\begin{figure}[H]
    \centering
    % \includegraphics[width=1.0\textwidth]{screenshot_home.png}
    \caption{Landing Page: visualizzazione pubblica con call-to-action per la registrazione.}
    \label{fig:home}
\end{figure}

\subsection{Sistema di Autenticazione e Registrazione}
La gestione degli accessi rappresenta uno snodo critico nel percorso utente. Il modulo di registrazione introduce una scelta logica dirimente sin dal primo contatto: la selezione del ruolo operativo.
In contrasto con i pattern di registrazione generici, l'utente è chiamato a dichiarare esplicitamente se intende fruire del servizio (\textit{Fruitore}) o far parte del network organizzativo (\textit{Volontario}). Questa profilazione precoce determina in modo permanente l'architettura dell'informazione e il layout successivo, personalizzando l'esperienza utente.

Il modulo richiede l'inserimento di dati obbligatori quali:
\begin{itemize}
    \item Nome e Cognome (per l'identificazione nelle liste presenze).
    \item Indirizzo Email (univoco, utilizzato come username).
    \item Password (con criteri minimi di sicurezza).
\end{itemize}

\begin{figure}[H]
    \centering
    % \includegraphics[width=0.7\textwidth]{screenshot_register.png}
    \caption{Modulo di registrazione con selettore per la tipologia di utenza (Fruitore/Volontario).}
    \label{fig:login_register}
\end{figure}

\subsection{Modulo di Richiesta Visita (Area Fruitore)}
A seguito dell'autenticazione, il \textit{Fruitore} interagisce con un'esplicita Call-to-Action che sblocca l'accesso all'ambiente dedicato alla schedulazione delle visite. Da un punto di vista dell'interazione, tale vista si discosta dal classico paradigma del carrello e-commerce in favore di un modulo guidato procedurale per l'inoltro di una "richiesta di servizio". 
Tale approccio minimizza il carico cognitivo dell'operatore, vincolando l'attenzione su input sequenziali mirati e semanticamente raggruppati.

I campi principali, validati per prevenire errori accidentali di sottomissione, richiedono:
\begin{itemize}
    \item \textbf{Dettagli Logistici:} Data preferita per la visita e numero di partecipanti (informazione cruciale per valutare la sostenibilità e il rapporto guida/visitatori).
    \item \textbf{Preferenze Context-Aware:} Selezione della lingua e campo testuale libero per documentare bisogni specifici del gruppo (es. accessibilità per disabilità motorie), promuovendo un design inclusivo.
\end{itemize}
Al completamento del flusso, il sistema fornisce un feedback contestuale di presa in carico, ponendo la pratica nello stato di sospeso in attesa di validazione asincrona da parte dell'amministrazione.

\begin{figure}[H]
    \centering
    % \includegraphics[width=0.8\textwidth]{screenshot_booking.png}
    \caption{Interfaccia di inserimento richiesta visita per l'utente Fruitore.}
    \label{fig:booking}
\end{figure}

\subsection{Gestione Disponibilità (Area Volontario)}
Il \textit{Volontario} ha come task principale l'aggiornamento della propria dashboard operativa. La vista adotta un'interfaccia a calendario (un solido modello mentale noto all'utente tramite applicativi esterni) che consente una mappatura diretta e manipolabile (\textit{Direct Manipulation}) delle proprie ore libere.
L'obiettivo cognitivo di questo cruscotto è agevolare l'inserimento dei dati essenziali per il futuro algoritmo organizzativo, riducendo le frizioni e promuovendo l'efficienza.

L'interazione prevede:
\begin{enumerate}
    \item Visualizzazione del mese corrente con evidenza dei giorni già impegnati.
    \item Possibilità di selezionare un giorno specifico e definire una fascia oraria (mattina/pomeriggio o orari puntuali).
    \item Feedback visivo immediato sulle disponibilità confermate.
\end{enumerate}

\begin{figure}[H]
    \centering
    % \includegraphics[width=0.9\textwidth]{screenshot_availability.png}
    \caption{Calendario interattivo per la gestione delle disponibilità del Volontario.}
    \label{fig:availability}
\end{figure}

\subsection{Pannello di Configurazione (Area Configuratore)}
Il ruolo direzionale (Configuratore) opera all'interno di un'area di back-office dedicata alla gestione infrastrutturale. Questa sezione è ideata come un hub di controllo per la componente statica dell'offerta turistica.
L'interazione è modellata su griglie dati (pattern CRUD - Create, Read, Update, Delete) pensate per le performance e l'accesso rapido all'informazione:
\begin{itemize}
    \item \textbf{Luoghi (Places):} Inserimento di nuovi siti visitabili, con dettagli descrittivi.
    \item \textbf{Tipologie di Visita:} Definizione delle categorie di tour offerti (es. "Visita Storica", "Percorso Naturalistico").
\end{itemize}

\begin{figure}[H]
    \centering
    % \includegraphics[width=1.0\textwidth]{screenshot_config_panel.png}
    \caption{Pannello di amministrazione per la gestione dei Luoghi e delle Tipologie di Visita.}
    \label{fig:config_panel}
\end{figure}

\subsection{Dashboard di Pianificazione (Visit Planning)}
La plancia direzionale di Visit Planning costituisce il nodo focale dell'applicativo, l'ambiente dove si attua la confluenza dei macro-processi generati da Fruitori e Volontari.
Data l'elevata densità di dati, l'interfaccia utente è calibrata specificamente per supportare il \textit{decision making} dell'amministratore, mitigando l'overload informativo.

La pagina è suddivisa in aree logiche:
\begin{itemize}
    \item \textbf{Lista Richieste Pendenti:} Un elenco delle visite richieste dai fruitori che necessitano di approvazione.
    \item \textbf{Strumenti di Assegnazione:} Per ogni richiesta, il sistema mostra i dettagli e permette di selezionare, tramite menu a tendina dinamici, un Luogo specifico e un Volontario compatibile (filtrato in base alla disponibilità temporale inserita).
    \item \textbf{Stato Avanzamento:} Indicatori che mostrano lo stato della visita (Richiesta, Confermata, Completata).
\end{itemize}
È in questa sezione che avviene il salvataggio definitivo dell'associazione Visita-Luogo-Volontario, scatenando eventuali notifiche di conferma.

\begin{figure}[H]
    \centering
    % \includegraphics[width=1.0\textwidth]{screenshot_visit_planning.png}
    \caption{Interfaccia operativa di Visit Planning: assegnazione risorse e gestione stati.}
    \label{fig:visit_planning}
\end{figure}
